\chapter{Giới thiệu doanh nghiệp và vị trí thực hiện}
\label{chap:gioi-thieu-don-vi}

\section{Giới thiệu doanh nghiệp}

\subsection{Thông tin chung về đơn vị}
Trung tâm Thông tin -- Thư viện UTH là đơn vị thuộc Trường Đại học Giao thông vận tải Thành phố Hồ Chí Minh, thực hiện các hoạt động thông tin -- thư viện và quản lý hệ thống công nghệ thông tin (CNTT) của Nhà trường. Trung tâm phục vụ nhu cầu học tập, giảng dạy, nghiên cứu khoa học và quản trị thông qua nguồn học liệu, dịch vụ thư viện và các hệ thống thông tin dùng chung~\cite{lic-gioi-thieu}.

Nhóm thực hiện Học kỳ doanh nghiệp (HKDN) tại Trung tâm với đề tài D7 -- Chatbot thư viện, dưới sự hướng dẫn của TS. Đặng Nhân Cách. Đề tài thuộc nhóm ứng dụng trí tuệ nhân tạo (AI) trong nâng cấp thư viện thông minh, tập trung hỗ trợ tra cứu, cung cấp thông tin giờ mở cửa, hỗ trợ gia hạn tài liệu và hướng dẫn sử dụng cơ sở dữ liệu điện tử~\cite{hkdn-de-tai}.

Thông tin khái quát về đơn vị được trình bày tại Bảng~\ref{tab:don-vi}.
\begin{table}[htbp]
\centering
\caption{Thông tin khái quát về đơn vị thực hiện}
\label{tab:don-vi}
\begin{tabularx}{\textwidth}{@{}P{3.7cm}Y@{}}
\toprule
\textbf{Nội dung} & \textbf{Thông tin}\\
\midrule
Tên đơn vị & Trung tâm Thông tin -- Thư viện UTH.\\
Đơn vị chủ quản & Trường Đại học Giao thông vận tải Thành phố Hồ Chí Minh.\\
Website & \url{https://lic.ut.edu.vn/}.\\
Lĩnh vực hoạt động & Thông tin -- thư viện; hệ thống và dịch vụ CNTT.\\
Đối tượng phục vụ & Cán bộ, giảng viên, sinh viên và người sử dụng các dịch vụ của Nhà trường.\\
\bottomrule
\end{tabularx}
\par\medskip
{\raggedright Nguồn: Tổng hợp từ thông tin giới thiệu của Trung tâm~\cite{lic-gioi-thieu}.\par}
\end{table}

\subsection{Quá trình hình thành và phát triển}
Hoạt động thư viện của Trường có quá trình phát triển từ năm 1991. Cùng với những thay đổi về tổ chức và nhu cầu phục vụ đào tạo, đơn vị từng bước mở rộng hoạt động từ cung cấp tài liệu sang kết hợp dịch vụ thư viện với quản lý dữ liệu và hạ tầng CNTT. Các mốc phát triển chính được tổng hợp tại Bảng~\ref{tab:lich-su} theo thông tin lịch sử do Trung tâm công bố~\cite{lic-lich-su}.

\begin{table}[htbp]
\centering
\caption{Các mốc hình thành và phát triển của Trung tâm}
\label{tab:lich-su}
\begin{tabularx}{\textwidth}{@{}P{1.4cm}Y@{}}
\toprule
\textbf{Năm} & \textbf{Nội dung}\\
\midrule
1991 & Thành lập Thư viện trực thuộc Phòng Đào tạo của Phân hiệu Đại học Hàng hải.\\
1998 & Tổ Thư viện trực thuộc Phòng Khoa học -- Công nghệ.\\
2001 & Tổ Thư viện trực thuộc Trường theo Quyết định số 41/TCCB.\\
2011 & Đổi tên Tổ Thư viện thành Thư viện Trường theo Quyết định số 480/QĐ/ĐH GTVT.\\
2014 & Thành lập Trung tâm Dữ liệu và Công nghệ thông tin theo Quyết định số 235/QĐ/ĐH GTVT.\\
2021 & Sáp nhập Thư viện và Trung tâm Dữ liệu và CNTT thành Trung tâm Thông tin -- Thư viện theo Quyết định số 856/QĐ-ĐHGTVT ngày 09/12/2021.\\
\bottomrule
\end{tabularx}
\par\medskip
{\raggedright Nguồn: Trung tâm Thông tin -- Thư viện UTH, “Lịch sử hình thành”~\cite{lic-lich-su}.\par}
\end{table}

Việc hợp nhất hai đơn vị tạo nên phạm vi hoạt động kết hợp giữa tổ chức, khai thác nguồn học liệu và xây dựng, vận hành hệ thống thông tin. Đây là đặc điểm của đơn vị tiếp nhận có liên quan trực tiếp đến hướng thực hiện đề tài chatbot hỗ trợ dịch vụ thư viện~\cite{lic-gioi-thieu}.

\subsection{Địa điểm hoạt động}
Trung tâm giới thiệu thư viện tại ba cơ sở của Nhà trường. Theo thông tin liên hệ trên website LIC, địa chỉ các cơ sở được trình bày trong Bảng~\ref{tab:co-so}~\cite{lic-trang-chu}.

\begin{table}[htbp]
\centering
\caption{Địa chỉ các cơ sở được công bố trên website LIC}
\label{tab:co-so}
\begin{tabularx}{\textwidth}{@{}P{1.5cm}Y@{}}
\toprule
\textbf{Cơ sở} & \textbf{Địa chỉ}\\
\midrule
1 & Số 2 Võ Oanh, phường Thạnh Mỹ Tây, TP. Hồ Chí Minh.\\
2 & Số 10 đường số 12, phường An Khánh, TP. Hồ Chí Minh.\\
3 & Số 70 Tô Ký, phường Trung Mỹ Tây, TP. Hồ Chí Minh.\\
\bottomrule
\end{tabularx}
\par\medskip
{\raggedright Nguồn: Trang chủ Trung tâm Thông tin -- Thư viện UTH~\cite{lic-trang-chu}.\par}
\end{table}

\section{Lĩnh vực hoạt động}

\subsection{Thông tin -- thư viện và phát triển học liệu}
Trung tâm thực hiện thu thập, xử lý, lưu trữ và cung cấp tài liệu phục vụ hoạt động đào tạo, học tập và nghiên cứu khoa học. Nguồn học liệu bao gồm tài liệu in, tài liệu điện tử và các cơ sở dữ liệu. Các hoạt động này gắn với nghiệp vụ tổ chức nguồn tin, hỗ trợ bạn đọc và khai thác tài nguyên thư viện~\cite{hkdn-de-tai}.

Các đầu mối tài nguyên được LIC giới thiệu gồm mục lục tra cứu trực tuyến (OPAC), Ebook UT--XD, ProQuest và mạng liên kết nguồn lực thông tin STINET. Những hệ thống này hỗ trợ người dùng tiếp cận các loại tài liệu và cơ sở dữ liệu theo điều kiện sử dụng của từng nguồn~\cite{lic-tai-nguyen}.

Thư viện số UTH tổ chức các nhóm tài nguyên như giáo trình -- bài giảng, khóa luận -- đồ án, luận văn và luận án tiến sĩ. Giao diện công khai cung cấp chức năng tìm kiếm và duyệt các nhóm tài nguyên, là một trong những đầu mối để người học tra cứu học liệu của Trường~\cite{uth-thu-vien-so}.

Đối với đề tài D7, thông tin nghiệp vụ và tài liệu hướng dẫn sử dụng dịch vụ là cơ sở xây dựng nội dung trả lời của chatbot. Nhóm tập trung vào những nội dung được đơn vị cho phép sử dụng, gắn câu trả lời với nguồn tham chiếu và phạm vi phục vụ cụ thể.

\subsection{Hệ thống thông tin và hạ tầng CNTT}
Trung tâm có chức năng tham mưu, xây dựng, quản lý và vận hành hệ thống CNTT toàn trường, phục vụ công tác quản trị, tuyển sinh, đào tạo và nghiên cứu~\cite{lic-gioi-thieu}. Các nhóm hạ tầng được nêu trong tài liệu HKDN gồm mạng nội bộ, Internet, Wi-Fi, máy chủ, dịch vụ web và hệ thống dữ liệu dùng chung~\cite{hkdn-de-tai}.

Đối với nhóm thực hiện D7, nội dung liên quan đến CNTT tập trung vào ứng dụng hỗ trợ người dùng, kết nối nguồn thông tin và điều kiện vận hành chatbot. Việc tiếp cận dữ liệu hoặc tích hợp với hệ thống thư viện tuân theo quyền truy cập và yêu cầu kỹ thuật của đơn vị.

\subsection{Dịch vụ hỗ trợ đào tạo và nghiên cứu}
Bên cạnh hoạt động thư viện và CNTT, Trung tâm công bố dịch vụ hỗ trợ xuất bản kỷ yếu, giáo trình đã nghiệm thu theo yêu cầu của tác giả và quy trình của nhà xuất bản~\cite{lic-xuat-ban}. Dịch vụ kiểm tra trùng lặp bằng Turnitin được giới thiệu cho các dạng công trình học thuật như luận văn, khóa luận, đồ án và bài báo~\cite{lic-trung-lap}.

Trung tâm duy trì các kênh hỗ trợ người sử dụng, trong đó có cổng hỗ trợ trực tuyến~\cite{lic-ho-tro}. Những dịch vụ này tạo bối cảnh nghiệp vụ để nhóm khảo sát nhu cầu hỏi đáp và xác định các nội dung phù hợp với phạm vi chatbot.

\section{Bộ phận tham gia}

\subsection{Phạm vi công việc tại Trung tâm}
Nhóm tham gia công việc ứng dụng CNTT hỗ trợ dịch vụ thư viện tại Trung tâm Thông tin -- Thư viện UTH. Đề tài D7 thuộc nhóm “Nâng cấp thư viện thông minh với AI” trong danh mục HKDN, với bốn nội dung nghiệp vụ là tra cứu, giờ mở cửa, gia hạn tài liệu và hướng dẫn sử dụng cơ sở dữ liệu điện tử~\cite{hkdn-de-tai}.

Phạm vi công việc của nhóm có sự liên hệ giữa nghiệp vụ thư viện và phát triển ứng dụng. Nội dung nghiệp vụ cung cấp căn cứ cho câu trả lời; phần ứng dụng tổ chức việc tiếp nhận câu hỏi, truy xuất thông tin và hiển thị kết quả cho người sử dụng. Sự phối hợp giữa hai mảng là cơ sở để xây dựng sản phẩm phù hợp với dịch vụ của đơn vị.

\subsection{Cơ chế phối hợp}
Nhóm trao đổi với Mentor về yêu cầu, phạm vi nhiệm vụ và tiến độ thực hiện. Các vấn đề về quy định thư viện, lịch phục vụ, hướng dẫn tài nguyên và đáp án kiểm thử cần được xác nhận qua đầu mối nghiệp vụ do đơn vị hướng dẫn. Các nội dung về dữ liệu, kết nối và triển khai được thực hiện theo điều kiện kỹ thuật và quyền hạn được cấp.

Theo hướng dẫn HKDN, Mentor là người trực tiếp hướng dẫn tại đơn vị, còn giảng viên đầu mối hỗ trợ kết nối giữa Nhà trường và đơn vị tiếp nhận~\cite{hkdn-huong-dan}. Nhóm thực hiện báo cáo tiến độ, tiếp nhận phản hồi và điều chỉnh công việc theo sự hướng dẫn của các đầu mối phụ trách.

\section{Vị trí và vai trò của sinh viên}

\subsection{Vị trí thực hiện}
Nhóm đảm nhiệm công việc khảo sát, phân tích yêu cầu và phát triển ứng dụng chatbot hỗ trợ dịch vụ thư viện. Hoạt động chuyên môn của sinh viên bao gồm tìm hiểu nghiệp vụ, chuẩn bị dữ liệu được phép sử dụng, thiết kế giải pháp, phát triển, kiểm thử và xây dựng tài liệu.

Trong giai đoạn đầu, trọng tâm công việc là khảo sát nguồn thông tin, xác định nhu cầu hỗ trợ bạn đọc và hoàn thiện yêu cầu của đề tài. Các chức năng tích hợp trực tiếp với hệ thống thư viện được xác định theo phạm vi nhiệm vụ và khả năng kết nối do đơn vị cung cấp.

\subsection{Nhiệm vụ chuyên môn}
Về khảo sát và phân tích, nhóm tìm hiểu các câu hỏi thường gặp, quy trình sử dụng dịch vụ và nguồn tài liệu chính thức. Kết quả của giai đoạn này là cơ sở xác định chức năng, dữ liệu đầu vào và tiêu chí đánh giá chatbot.

Về thiết kế và phát triển, nhóm xây dựng luồng hội thoại, tổ chức việc truy xuất thông tin và trình bày câu trả lời có nguồn tham chiếu. Các nội dung tra cứu, lịch phục vụ, hướng dẫn gia hạn và hướng dẫn cơ sở dữ liệu được triển khai theo yêu cầu đã thống nhất với đơn vị.

Về kiểm thử và bàn giao, nhóm chuẩn bị tình huống nghiệp vụ, đối chiếu kết quả với tài liệu được duyệt, ghi nhận lỗi và hoàn thiện hướng dẫn sử dụng, vận hành. Những nhiệm vụ này gắn với yêu cầu đầu ra của đề tài HKDN, bao gồm mã nguồn, tài liệu kỹ thuật, báo cáo, demo và bảng đánh giá độ chính xác đối với sản phẩm AI~\cite{hkdn-de-tai}.

\subsection{Trách nhiệm trong quá trình thực hiện}
Nhóm tuân thủ nội quy, thời gian làm việc và sự phân công của đơn vị; chủ động trao đổi khi phát sinh khó khăn. Việc sử dụng dữ liệu, mã nguồn và tài liệu nội bộ tuân theo yêu cầu bảo mật; các thông tin này không được tự ý cung cấp hoặc công bố~\cite{hkdn-huong-dan}.

Báo cáo tiến độ thể hiện công việc đã thực hiện, kết quả đạt được, khó khăn, công việc đang làm và kế hoạch tiếp theo. Hình thức và tần suất báo cáo thực hiện theo quy định của Mentor hoặc đơn vị~\cite{hkdn-huong-dan}. Mỗi thành viên cập nhật nhật ký công việc, lưu minh chứng và chịu trách nhiệm đối với phần nội dung mình phụ trách.

\section{Mentor và nhóm làm việc}

\subsection{Mentor tại đơn vị}
Mentor hướng dẫn nhóm là \textbf{TS. Đặng Nhân Cách}. Theo thông tin nhân sự được LIC công bố, TS. Đặng Nhân Cách giữ chức vụ Giám đốc Trung tâm Thông tin -- Thư viện~\cite{lic-gioi-thieu}.

Mentor định hướng công việc, hướng dẫn quy trình, theo dõi tiến độ, góp ý và đánh giá kết quả thực hiện của sinh viên theo hướng dẫn HKDN~\cite{hkdn-huong-dan}. Với đề tài D7, nhóm trao đổi các nội dung về phạm vi chức năng, nguồn dữ liệu, yêu cầu nghiệp vụ và các mốc kiểm tra để tổ chức công việc phù hợp với điều kiện của đơn vị.

\subsection{Thành viên nhóm}
Nhóm thực hiện gồm bốn sinh viên, được trình bày trong Bảng~\ref{tab:thanh-vien}. Thời hạn kết thúc HKDN của nhóm là ngày \textbf{28/12/2026}.

\begin{table}[htbp]
\centering
\caption{Danh sách thành viên nhóm thực hiện}
\label{tab:thanh-vien}
\begin{tabularx}{\textwidth}{@{}cP{3.1cm}Y@{}}
\toprule
\textbf{STT} & \textbf{Mã số sinh viên} & \textbf{Họ và tên}\\
\midrule
1 & 075205002473 & Lê Sĩ Trung\\
2 & 070205005207 & Nguyễn Việt Dũng\\
3 & 075205014655 & Hoàng Lê Ngọc Huy\\
4 & 075205015362 & Hoàng Đặng Văn Kỳ\\
\bottomrule
\end{tabularx}
\par\medskip
{\raggedright Nguồn: Thông tin nhóm thực hiện.\par}
\end{table}

\subsection{Tổ chức công việc của nhóm}
Nhóm dự kiến tổ chức công việc theo bốn mảng: nghiệp vụ và dữ liệu; xử lý phía máy chủ và AI; giao diện và tích hợp; kiểm thử và tài liệu bàn giao. Cách tổ chức này phục vụ việc phân chia nhiệm vụ và phối hợp giữa các thành viên trong quá trình thực hiện D7.

Trong mảng nghiệp vụ và dữ liệu, công việc tập trung vào tổng hợp yêu cầu, lựa chọn nguồn và chuẩn bị nội dung trả lời. Mảng xử lý phía máy chủ và AI phụ trách tiếp nhận câu hỏi, truy xuất thông tin và tạo phản hồi. Mảng giao diện và tích hợp tổ chức tương tác với người dùng, hiển thị nguồn và kết quả tra cứu. Mảng kiểm thử và tài liệu chuẩn bị tình huống đánh giá, ghi nhận kết quả và tổng hợp hướng dẫn bàn giao.

% Ban ghi bien soan va cac thong tin can bo sung nam tai ghi-chu-bien-tap.tex.
